% Options for packages loaded elsewhere
% Options for packages loaded elsewhere
\PassOptionsToPackage{unicode}{hyperref}
\PassOptionsToPackage{hyphens}{url}
%
\documentclass[
  11pt,
  ignorenonframetext,
  aspectratio=169]{beamer}
\newif\ifbibliography
\usepackage{pgfpages}
\setbeamertemplate{caption}[numbered]
\setbeamertemplate{caption label separator}{: }
\setbeamercolor{caption name}{fg=normal text.fg}
\beamertemplatenavigationsymbolsempty
% remove section numbering
\setbeamertemplate{part page}{
  \centering
  \begin{beamercolorbox}[sep=16pt,center]{part title}
    \usebeamerfont{part title}\insertpart\par
  \end{beamercolorbox}
}
\setbeamertemplate{section page}{
  \centering
  \begin{beamercolorbox}[sep=12pt,center]{section title}
    \usebeamerfont{section title}\insertsection\par
  \end{beamercolorbox}
}
\setbeamertemplate{subsection page}{
  \centering
  \begin{beamercolorbox}[sep=8pt,center]{subsection title}
    \usebeamerfont{subsection title}\insertsubsection\par
  \end{beamercolorbox}
}
% Prevent slide breaks in the middle of a paragraph
\widowpenalties 1 10000
\raggedbottom
\AtBeginPart{
  \frame{\partpage}
}
\AtBeginSection{
  \ifbibliography
  \else
    \frame{\sectionpage}
  \fi
}
\AtBeginSubsection{
  \frame{\subsectionpage}
}
\usepackage{iftex}
\ifPDFTeX
  \usepackage[T1]{fontenc}
  \usepackage[utf8]{inputenc}
  \usepackage{textcomp} % provide euro and other symbols
\else % if luatex or xetex
  \usepackage{unicode-math} % this also loads fontspec
  \defaultfontfeatures{Scale=MatchLowercase}
  \defaultfontfeatures[\rmfamily]{Ligatures=TeX,Scale=1}
\fi
\usepackage{lmodern}

\ifPDFTeX\else
  % xetex/luatex font selection
\fi
% Use upquote if available, for straight quotes in verbatim environments
\IfFileExists{upquote.sty}{\usepackage{upquote}}{}
\IfFileExists{microtype.sty}{% use microtype if available
  \usepackage[]{microtype}
  \UseMicrotypeSet[protrusion]{basicmath} % disable protrusion for tt fonts
}{}
\makeatletter
\@ifundefined{KOMAClassName}{% if non-KOMA class
  \IfFileExists{parskip.sty}{%
    \usepackage{parskip}
  }{% else
    \setlength{\parindent}{0pt}
    \setlength{\parskip}{6pt plus 2pt minus 1pt}}
}{% if KOMA class
  \KOMAoptions{parskip=half}}
\makeatother


\usepackage{longtable,booktabs,array}
\usepackage{calc} % for calculating minipage widths
\usepackage{caption}
% Make caption package work with longtable
\makeatletter
\def\fnum@table{\tablename~\thetable}
\makeatother
\usepackage{graphicx}
\makeatletter
\newsavebox\pandoc@box
\newcommand*\pandocbounded[1]{% scales image to fit in text height/width
  \sbox\pandoc@box{#1}%
  \Gscale@div\@tempa{\textheight}{\dimexpr\ht\pandoc@box+\dp\pandoc@box\relax}%
  \Gscale@div\@tempb{\linewidth}{\wd\pandoc@box}%
  \ifdim\@tempb\p@<\@tempa\p@\let\@tempa\@tempb\fi% select the smaller of both
  \ifdim\@tempa\p@<\p@\scalebox{\@tempa}{\usebox\pandoc@box}%
  \else\usebox{\pandoc@box}%
  \fi%
}
% Set default figure placement to htbp
\def\fps@figure{htbp}
\makeatother





\setlength{\emergencystretch}{3em} % prevent overfull lines

\providecommand{\tightlist}{%
  \setlength{\itemsep}{0pt}\setlength{\parskip}{0pt}}



 


%% BEAMER THEME ITHACA 2013: Main tex file for compiling
%$ Compile this file. 
%%
%% Copyright 2013 by Raul Gomez
%% This file may be distributed and/or modified
%% 	1. under the LaTeX Project Public License and/or
%% 	2. under the GNU Public License.
%% 
%% This theme is based on Flip's Theme by Flip Tanedo.

%% Please see notes.txt for comments on Beamer Theme Ithaca 2013
%% By default, this template is meant to be run with XeLaTeX (for fonts)
%% To run in PDFLaTeX, remove fontspec and any font commands

%% Discussion of Beamer vs XeLaTeX vs LuaLaTeX
%% http://tex.stackexchange.com/questions/29497/xelatex-preventing-beamer-from-using-different-backgrounds



%\documentclass[12 pt,hyperref={pdfpagelabels=false}]{beamer}
%\documentclass[aspectratio=169]{beamer}



\usetheme[
%	bullet=square,           % Default option: circle
%	nobigpagenumber,         % No big circled page number on lower right
%       nofootline,		 % No footline
	grey                    % For a neutral grey background color
%	watermark=BG_both,	 % png file for the watermark
	]{Ithaca}



%%%%%%%%%%
% FONTS %
%%%%%%%%%%

%% Default font: lmodern, doesn't require fontspec % solves some default warnings
\usepackage[T1]{fontenc}
\usepackage{lmodern}	
\usepackage{mathptmx}
\usepackage{frutiger}		

%% Protects fonts from Beamer screwing with them
%% http://tex.stackexchange.com/questions/10488/force-computer-modern-in-math-mode
\usefonttheme{professionalfonts}



%%%%%%%%%%%%%%%%%%%%%%%%
% Usual LaTeX Packages %
%%%%%%%%%%%%%%%%%%%%%%%%

\usepackage{amsmath}
\usepackage{amsfonts}
\usepackage{amssymb}
\usepackage{graphicx}
\usepackage{mathrsfs} 			% For Weinberg-esque letters
\usepackage{cancel}			% For "SUSY-breaking" symbol
\usepackage{slashed}   		        % for slashed characters in math mode
\usepackage{bbm}                	% for \mathbbm{1} (unit matrix)
\usepackage{amsthm}			% For theorem environment
\usepackage{multirow}			% For multi row cells in table
\usepackage{young}			% For Young Tableaux
\usepackage{xspace}			% For spacing after commands
\usepackage{wrapfig}			% for Text wrap around figures
\usepackage{framed}
\usepackage{pgf}
\usepackage[table]{xcolor}
\usepackage{tikz}
\usepackage{hyperref}
\usepackage{float}



\graphicspath{{images/}}		% Put all images in this directory. Avoids clutter.

\newcommand{\cG}{\mathcal{T}}
\newcommand{\supp}{\operatorname{supp}}
\newcommand{\Span}{\operatorname{Span}}
\newcommand{\R}{\mathbb{R}}
\newcommand{\C}{\mathbb{C}}
\newcommand{\N}{\mathbb{N}}
\newcommand{\Z}{\mathbb{Z}}
\newcommand{\RP}{\mathbb{RP}}
\newcommand{\g}{\mathfrak{g}}
\newcommand{\n}{\mathfrak{n}}
\newcommand{\m}{\mathfrak{m}}
\newcommand{\p}{\mathfrak{p}}
\renewcommand{\sl}{\mathfrak{sl}}
\renewcommand{\S}{\mathscr{S}}
\newcommand{\Hom}{\operatorname{Hom}}
\newcommand{\SL}{\operatorname{SL}}
\newcommand{\PGL}{\operatorname{PGL}}
\renewcommand{\a}{\mathfrak{a}}
\newcommand{\ind}{\operatorname{ind}}
\newcommand{\Ext}{\operatorname{Ext}}
\newcommand{\Tor}{\operatorname{Tor}}
\newcommand{\hot}{\widehat{\otimes}}
\newcommand{\Fre}{{Fr\'{e}chet \,}}
\newcommand{\Sc}{{\mathcal S}}
\newcommand{\sgn}{\operatorname{sgn}}
\newcommand{\Orb}{\mathcal{O}}
\newcommand{\Cov}{\operatorname{Cov}}
\newcommand{\Var}{\operatorname{Var}}
\newcommand{\Tr}{\operatorname{Tr}}
\renewcommand{\N}{\mathcal{N}}
\newcommand{\diag}{\operatorname{diag}}
\newcommand{\E}{\mathbb{E}}


%%%%%%%%%%%%%%%%%%%%%%%%%%%%%%%%%%%%%%%%%%%%%%
% Theorem, Definition and Remark Environments%
%%%%%%%%%%%%%%%%%%%%%%%%%%%%%%%%%%%%%%%%%%%%%%

\theoremstyle{plain}
%\newtheorem{proposition}[theorem]{Proposition}
%\newtheorem{conjecture}[theorem]{Conjecture}
%\newtheorem{claim}[theorem]{Claim}
%\newtheorem{algorithm}[theorem]{Algorithm}
%\newtheorem{axiom}[theorem]{Axiom}
%\newtheorem{criterion}[theorem]{Criterion}

\theoremstyle{definition}
%\newtheorem{condition}[theorem]{Condition}
%\newtheorem{exercise}[theorem]{Exercise}
%\newtheorem{notation}[theorem]{Notation}
%\newtheorem{case}[theorem]{Case}
%\newtheorem{question}[theorem]{Question}

\theoremstyle{remark}
%\newtheorem{remark}[theorem]{Remark}
%\newtheorem{remarks}[theorem]{Remarks}
%\newtheorem{summary}[theorem]{Summary}
%\newtheorem{observation}[theorem]{Observation}
%\newtheorem{conclusion}[theorem]{Conclusion}
%\newtheorem{acknowledgement}[theorem]{Acknowledgement}


\newcounter{saveenumi}
\newcommand{\seti}{\setcounter{saveenumi}{\value{enumi}}}
\newcommand{\conti}{\setcounter{enumi}{\value{saveenumi}}}
%%%%%%%%%%%%%%%%%%%%%%%%%%%%%%%%%%%%%%
% Section and subsection adjustments %
%%%%%%%%%%%%%%%%%%%%%%%%%%%%%%%%%%%%%%
\AtBeginSection{}
\AtBeginSubsection{}	 




%%%%%%%%%%%%
% Preamble %
%%%%%%%%%%%%

\title[MA2003B]{8 Análisis factorial: métodos prácticos}
\subtitle{MA2003B Aplicación de métodos multivariados \\ en ciencia de datos}
\institute[Tec de Monterrey]{EIC}
\date{Campus Monterrey}
\logo{
\includegraphics[height=1.0cm]{tecnologico-de-monterrey-blue}}
\makeatletter
\@ifpackageloaded{caption}{}{\usepackage{caption}}
\AtBeginDocument{%
\ifdefined\contentsname
  \renewcommand*\contentsname{Table of contents}
\else
  \newcommand\contentsname{Table of contents}
\fi
\ifdefined\listfigurename
  \renewcommand*\listfigurename{List of Figures}
\else
  \newcommand\listfigurename{List of Figures}
\fi
\ifdefined\listtablename
  \renewcommand*\listtablename{List of Tables}
\else
  \newcommand\listtablename{List of Tables}
\fi
\ifdefined\figurename
  \renewcommand*\figurename{Figure}
\else
  \newcommand\figurename{Figure}
\fi
\ifdefined\tablename
  \renewcommand*\tablename{Table}
\else
  \newcommand\tablename{Table}
\fi
}
\@ifpackageloaded{float}{}{\usepackage{float}}
\floatstyle{ruled}
\@ifundefined{c@chapter}{\newfloat{codelisting}{h}{lop}}{\newfloat{codelisting}{h}{lop}[chapter]}
\floatname{codelisting}{Listing}
\newcommand*\listoflistings{\listof{codelisting}{List of Listings}}
\makeatother
\makeatletter
\makeatother
\makeatletter
\@ifpackageloaded{caption}{}{\usepackage{caption}}
\@ifpackageloaded{subcaption}{}{\usepackage{subcaption}}
\makeatother

\usepackage{bookmark}
\IfFileExists{xurl.sty}{\usepackage{xurl}}{} % add URL line breaks if available
\urlstyle{same}
\hypersetup{
  pdfauthor={rgomez@tec.mx},
  hidelinks,
  pdfcreator={LaTeX via pandoc}}


\author{rgomez@tec.mx}
\date{}

\begin{document}


\begin{frame}[plain]{}
\phantomsection\label{section}
\addtocounter{framenumber}{-1}

\begin{center}
\Huge\textcolor{structure}{\inserttitle}

\vfill

\Large\insertsubtitle

\vfill

\vspace{1em}

\vfill


\large \raisebox{.25cm}{Raul Gomez} \qquad \insertlogo


\vfill

\end{center}
\end{frame}

\section{Introducción al Tidyverse}\label{introducciuxf3n-al-tidyverse}

\section{Los momentos estandarizados}\label{los-momentos-estandarizados}

\section{Pruebas de normalidad}\label{pruebas-de-normalidad}

\section{Introducción al análisis
multivariante}\label{introducciuxf3n-al-anuxe1lisis-multivariante}

\section{Distribución normal
multivariante}\label{distribuciuxf3n-normal-multivariante}

\subsection{Definición y propiedades
básicas}\label{definiciuxf3n-y-propiedades-buxe1sicas}

\begin{frame}{Definición y propiedades básicas}
\phantomsection\label{definiciuxf3n-y-propiedades-buxe1sicas-1}
\begin{definition}
Sea $\mu$ un vector (renglón) y $\Sigma$ una matriz simétrica definida positiva. Decimos que un vector aleatorio (renglón) $X$ tiene \alert{distribución normal multivariante} con media $\mu$ y matriz de covarianzas $\Sigma$, en símbolos $X \sim \mathcal{N}_{p}(\mu, \Sigma)$, si su densidad es
$$
\frac{1}{(2\pi)^{p/2} |\Sigma|^{1/2}} \exp\left(-\frac{1}{2}(X - \mu) \Sigma^{-1} (X - \mu)^{T}\right)dX
$$
\end{definition}
\end{frame}

\begin{frame}{}
\phantomsection\label{section-1}
\begin{remark} Algunas observaciones sobre la definición:
\begin{enumerate}
\item Al igual que en el caso univariado, la normal multivariante queda completamente determinada por $\mu$ y $\Sigma$. \pause
\item La matriz $\Sigma$ debe ser simétrica y definida positiva para que la distribución esté bien definida y el exponente sea negativo definido. \pause
\item El término 
$$
|\Sigma| = \det(\Sigma)
$$
es la \alert{varianza generalizada} de la distribución. Puede interpretarse como una medida de la "dispersión" en el espacio de dimensión $p$. \pause
\item El término $(X - \mu) \Sigma^{-1} (X - \mu)^{T}$ es el cuadrado de la distancia de Mahalanobis entre $X$ y la media $\mu$. 
\end{enumerate}
\end{remark}
\end{frame}

\begin{frame}{}
\phantomsection\label{section-2}
Como en el caso univariado, la normal multivariante es destacada por el
Teorema Central del Límite, que a grandes rasgos dice que la suma de
muchas variables i.i.d. (debidamente estandarizadas) converge a una
normal.
\end{frame}

\begin{frame}{Teorema central del límite (multivariante)}
\phantomsection\label{teorema-central-del-luxedmite-multivariante}
\begin{theorem}
Sea $\{X_i\}_{i=1}^{\infty}$ una sucesión de vectores aleatorios i.i.d. en $\mathbb{R}^p$ con media $\mu$ y covarianzas $\Sigma$. Defina la media muestral
\[
\bar{X}_n = \frac{1}{n}\sum_{i=1}^n X_i.
\]
Entonces, cuando $n \to \infty$,
\[
\sqrt{n}\,(\bar{X}_n - \mu) \;\xrightarrow{d}\; \mathcal{N}_p(0, \Sigma),
\]
donde $\xrightarrow{d}$ denota convergencia en distribución y $\mathcal{N}_p(0, \Sigma)$ es la normal multivariante de dimensión $p$ con media $0$ y covarianzas $\Sigma$.
\end{theorem}
\end{frame}

\begin{frame}{}
\phantomsection\label{section-3}
\begin{remark}
Aun cuando el TCL destaca a la normal multivariante como especialmente importante, no es la única distribución relevante en estadística multivariante.

\bigpause 

De hecho, la mayoría de los datos reales no siguen una normal, y muchos métodos están diseñados para datos no normales.

\bigpause

Además, métodos importantes como regresión lineal, ACP y clustering no producen buenos resultados cuando se aplican a una normal multivariante.
\end{remark}
\end{frame}

\subsection{Contornos de la normal
multivariante}\label{contornos-de-la-normal-multivariante}

\begin{frame}{Contornos de la normal multivariante}
\phantomsection\label{contornos-de-la-normal-multivariante-1}
Observe que la función de densidad de la normal multivariante es
gaussiana en \(p\) dimensiones.

\bigpause

Esto implica que los contornos son elipsoides centrados en \(\mu\) y
cuya forma está controlada por \(\Sigma\).

\bigpause

La probabilidad asociada a estos contornos está controlada por la
distribución \(\chi^2\) en el siguiente sentido:
\end{frame}

\begin{frame}{}
\phantomsection\label{section-4}
Sea \(0 \leq \varpi \leq 1\) un nivel de probabilidad, y sea \(c\) tal
que \[
\mathbb{P}(\chi^2_{p} \leq c) = \varpi
\]

\bigpause

Entonces, el contorno que contiene una proporción \(\varpi\) de la masa
de probabilidad está dado por los puntos \(X\) que satisfacen \[
(X - \mu) \Sigma^{-1} (X - \mu)^{T} = c 
\]
\end{frame}

\begin{frame}{}
\phantomsection\label{section-5}
En la siguiente figura mostramos los contornos de la normal bivariante
con distintas matrices de covarianzas.

\bigpause

Los contornos son elipses centradas en \(\mu = (0, 0)\), y su forma está
determinada por \(\Sigma\).

\bigpause

El gradiente de color representa la densidad: colores más oscuros
indican mayor densidad.
\end{frame}

\begin{frame}{}
\phantomsection\label{section-6}
\begin{figure}[H]

{\centering \pandocbounded{\includegraphics[keepaspectratio]{5_normal_multivariate_distribution_files/figure-beamer/unnamed-chunk-2-1.pdf}}

}

\caption{Distribución normal bivariante con distintas matrices de
covarianzas}

\end{figure}%
\end{frame}

\subsection{Estandarización de la normal
multivariante}\label{estandarizaciuxf3n-de-la-normal-multivariante}

\begin{frame}{Estandarización de la distribución normal multivariante}
\phantomsection\label{estandarizaciuxf3n-de-la-distribuciuxf3n-normal-multivariante}
Al igual que en el caso univariado, dado
\(X \sim \mathcal{N}_{p}(\mu, \Sigma)\), podemos estandarizarlo para
obtener \(Z \sim \mathcal{N}_{p}(0, I_p)\), donde \(I_p\) es la
identidad de tamaño \(p\).

\bigpause

Para describir este proceso, recordemos algunos resultados de álgebra
lineal.
\end{frame}

\begin{frame}{Singular Value Decomposition}
\phantomsection\label{singular-value-decomposition}
\begin{theorem}
Sea $\Sigma$ simétrica definida positiva de tamaño $p \times p$. Entonces existe una matriz ortogonal $Q$ y una diagonal $\Lambda$ con entradas positivas tales que 
$$
\Sigma = Q \Lambda Q^{T}
$$
\end{theorem}

\pause

\begin{remark}
Recuerde que $Q$ es ortogonal si $Q^{T}Q = I_p$, donde $I_p$ es la identidad de tamaño $p$. Sus columnas son ortonormales y forman una base de $\mathbb{R}^p$.
\end{remark}
\end{frame}

\begin{frame}{}
\phantomsection\label{section-7}
\begin{remark}
La descomposición en valores singulares (SVD) es una herramienta que descompone una matriz en partes constitutivas.

\bigpause 

Para $\Sigma$, la SVD la expresa en términos de sus \alert{valores singulares} (entradas diagonales de $\Lambda$, también llamados \alert{autovalores}) y sus \alert{vectores singulares} (columnas de $Q$, también llamados \alert{autovectores}).

\bigpause

Esta descomposición es muy útil para entender la geometría de la normal multivariante.
\end{remark}
\end{frame}

\begin{frame}{}
\phantomsection\label{section-8}
Obsérvese que, como \(\Sigma\) es definida positiva, las entradas
diagonales de \[
\Lambda = \left[
\begin{array}{cccc}
\lambda_1 & 0 & \cdots & 0 \\
0 & \lambda_2 & \cdots & 0 \\
\vdots & \vdots & \ddots & \vdots \\
0 & 0 & \cdots & \lambda_p
\end{array}
\right]
\] son todas positivas.

\bigpause

En particular, podemos definir la matriz \[
D = \left[
\begin{array}{cccc}
\sqrt{\lambda_1} & 0 & \cdots & 0 \\
0 & \sqrt{\lambda_2} & \cdots & 0 \\
\vdots & \vdots & \ddots & \vdots \\
0 & 0 & \cdots & \sqrt{\lambda_p}
\end{array}
\right]
\] que satisface \(D^2 = \Lambda\).
\end{frame}

\begin{frame}{}
\phantomsection\label{section-9}
Si ahora tomamos \[ 
\sigma = D Q^{T}
\]

\bigpause

entonces \[
\Sigma = \sigma^{T} \sigma
\]

\bigpause

En particular, podemos definir la \alert{estandarización} de \(X\) como
\[
Z = (X - \mu) \sigma^{-1} = (X - \mu) Q D^{-1} = (X - \mu) Q \Lambda^{-1/2}
\] in which case \(Z \sim \mathcal{N}_{p}(0, I_p)\).
\end{frame}

\subsection{Interpretación geométrica de la
SVD}\label{interpretaciuxf3n-geomuxe9trica-de-la-svd}

\begin{frame}{Interpretación geométrica de la SVD}
\phantomsection\label{interpretaciuxf3n-geomuxe9trica-de-la-svd-1}
Antes consideramos la normal multivariante con media \[ 
\mu = [0, 0] 
\] y matriz de covarianzas \[
\Sigma = \left[
\begin{array}{cc}
1 & 0.6 \\
0.6 & 1 
\end{array} 
\right] 
\]

\bigpause

Según la SVD, podemos escribir esta matriz de covarianzas como \[ 
\Sigma = Q \Lambda Q^{T} 
\] for some orthogonal matrix \(Q\) and diagonal matrix \(\Lambda\).
\end{frame}

\begin{frame}{}
\phantomsection\label{section-10}
En este caso obtenemos: \[
Q = \left[
\begin{array}{cc}
 0.7071068 & -0.7071068 \\
 0.7071068 & 0.7071068
\end{array}
\right]
\] y \[
\Lambda = \left[
\begin{array}{cc}
 1.6 & 0 \\
0 & 0.4
\end{array}
\right]
\] en donde los vectores singulares son las columnas de \(Q\) y los
valores singulares son las entradas diagonales de \(\Lambda\).

\bigpause

Geometrically, the singular vectors represent the directions of the axes
of the ellipsoid that defines the contours of the normal multivariate
distribution, and the square root of the singular values represent the
lengths of these axes.
\end{frame}

\begin{frame}{}
\phantomsection\label{section-11}
\begin{center}
\pandocbounded{\includegraphics[keepaspectratio]{5_normal_multivariate_distribution_files/figure-beamer/unnamed-chunk-4-1.pdf}}
\end{center}
\end{frame}

\subsection{Clausura bajo transformaciones
lineales}\label{clausura-bajo-transformaciones-lineales}

\begin{frame}{Clausura bajo transformaciones lineales}
\phantomsection\label{clausura-bajo-transformaciones-lineales-1}
\begin{theorem}
Sea $X \sim \mathcal{N}_{p}(\mu, \Sigma)$. Para cualquier matriz $A$ de tamaño $m \times p$ y vector $b$ de tamaño $m$, el vector aleatorio 
$$
Y = AX + b 
$$ 
también tiene distribución normal multivariante: $Y \sim \mathcal{N}_{m}(A\mu + b, A\Sigma A^{T})$. 
\end{theorem}
\end{frame}

\begin{frame}{}
\phantomsection\label{section-12}
\begin{remark} 
El teorema dice que si aplicamos una transformación lineal a un vector normal multivariante, el resultado es también normal multivariante.

\bigpause 

La media y la covarianza del transformado vienen dadas por las fórmulas del teorema.

\bigpause

Es útil en aplicaciones como regresión, donde transformamos datos preservando la normalidad de residuales.

\bigpause

En la práctica se usa para derivar la distribución de combinaciones lineales de variables aleatorias.

\end{remark}
\end{frame}

\begin{frame}{}
\phantomsection\label{section-13}
\begin{example} 
Sea $X \sim \mathcal{N}_{2}(\mu, \Sigma)$ con media $\mu = (1, 2)$ y matriz de covarianzas 
$$ 
\Sigma = \left[ 
\begin{array}{cc} 
1 & 0.5 \\ 
0.5 & 2 
\end{array} 
\right] 
$$ 

\pause

Ahora, considere la transformación lineal 
$$ 
Y = \begin{bmatrix} 2 & 1 \\ 0 & 3 \end{bmatrix} X + \begin{bmatrix} 0 \\ 1 \end{bmatrix} 
$$ 

\pause

En otras palabras, 
\begin{align*}
Y_1 & = 2X_1 + X_2 \\ 
Y_2 & = 3X_2 + 1 
\end{align*}
\end{example}
\end{frame}

\begin{frame}{}
\phantomsection\label{section-14}
\begin{example}

Usando la clausura bajo transformaciones lineales, obtenemos media y covarianzas de $Y$:
\begin{align*} 
\text{Mean: } & A\mu + b = \begin{bmatrix} 2 & 1 \\ 0 & 3 \end{bmatrix} \begin{bmatrix} 1 \\ 2 \end{bmatrix} + \begin{bmatrix} 0 \\ 1 \end{bmatrix} = \begin{bmatrix} 4 \\ 7 \end{bmatrix} \\ 
\text{Covariance: } & A\Sigma A^{T} = \begin{bmatrix} 2 & 1 \\ 0 & 3 \end{bmatrix} \begin{bmatrix} 1 & 0.5 \\ 0.5 & 2 \end{bmatrix} \begin{bmatrix} 2 & 0 \\ 1 & 3 \end{bmatrix} = \begin{bmatrix} 8 & 9 \\ 9 & 18 \end{bmatrix} 
\end{align*} 

\pause

Así,  
$$Y \sim \mathcal{N}_{2}\left(\begin{bmatrix} 4 \\ 7 \end{bmatrix}, \begin{bmatrix} 8 & 9 \\ 9 & 18 \end{bmatrix}\right)
$$ 
\end{example}
\end{frame}

\begin{frame}{}
\phantomsection\label{section-15}
\begin{remark} Nótese que proyectar a un subespacio es lineal; por tanto, cualquier subconjunto de componentes de una normal multivariante también es normal multivariante.

\bigpause

Esta propiedad permite analizar subconjuntos de variables manteniendo la suposición de normalidad. 
\end{remark}
\end{frame}

\subsection{Pruebas de normalidad
multivariante}\label{pruebas-de-normalidad-multivariante}

\begin{frame}{Pruebas de normalidad multivariante}
\phantomsection\label{pruebas-de-normalidad-multivariante-1}
Existen varias pruebas para verificar si un conjunto de datos sigue una
normal multivariante. Son esenciales porque muchos métodos asumen
normalidad. Algunas comunes y su implementación en R:

\begin{enumerate}
  \item \alert{Prueba de Mardia:}  
    Basada en asimetría y curtosis multivariantes.  
    Para normal multivariante, asimetría ≈ 0 y curtosis ≈ \(p(p+2)\). En R: \texttt{MVN::mvn(tus\_datos, mvn\_test = "mardia")}.
    
  \item \alert{Prueba de Henze–Zirkler (HZ):}  
    Usa la función característica empírica para medir desvíos de normalidad. Buen poder global; adecuada en dimensiones moderadas. En R: \texttt{MVN::mvn(tus\_datos, mvn\_test = "hz")}.
    
\seti 
\end{enumerate}
\end{frame}

\begin{frame}{}
\phantomsection\label{section-16}
\begin{enumerate} 

\conti

  \item \alert{Prueba de Royston:}  
    Extensión multivariante de Shapiro–Wilk (producto de estadísticas univariadas aplicado a componentes principales). Fácil de interpretar; pierde poder en alta dimensión. En R: \texttt{MVN::mvn(tus\_datos, mvn\_test = "royston")}.
    
  \item \alert{Prueba de Doornik–Hansen:}  
    Combina asimetría y curtosis en una prueba ómnibus, alternativa robusta a Mardia. En R: \texttt{MVN::mvn(tus\_datos, mvn\_test = "dh")}.
    
  \item \alert{Métodos gráficos (Q–Q de chi‑cuadrada):}  
    Graficar distancias de Mahalanobis contra cuantiles de $\chi^2$. Desviaciones de la recta indican no normalidad. En R: \texttt{multivariate\_diagnostic\_plot(data = tus\_datos, type = "qq")}.
\end{enumerate}
\end{frame}




\end{document}
